% !TeX root = ../../../../../main.tex
% chapters/sections/chapter5/subsections/a_shock/subsubsections/interpretations/fall.tex

\subsubsection{\(i_t^H\)が下落する政策群（CPLT, IT-CPI, TR-CPI）の解釈}
\label{sec:chapter5_a_shock_interpretations_fall}

\begin{enumerate}
    \item
    \label{itm:chapter5_a_shock_interpretations_fall_first}
        \textbf{1期：\(y_t^H,c_t^{H\to W}\)の下落}
        \begin{description}[style=nextline]
            \item[\(\hat{a}_1^H\)の負への下落と\(i_1^H\)の下落]
            \label{itm:chapter5_a_shock_interpretations_fall_first_a_H}
                \(a\)ショックが発生し、予期せず\(\hat{a}_1^H\)が負へ大きく下落する。
                負へ下落した\(\hat{a}_t^H\)は自己回帰過程の式
                \eqref{form:chapter4_shocks_a_eq}
                にしたがって増加し、やがて初期定常状態\(0\)へ回帰する。
                また後述の諸変数の変動により\(i_1^H\)は金融政策の式を通じて下落する。
                下落した\(i_t^H\)は金融政策の式にしたがって将来のいずれかの期から上昇しはじめ、
                やがて初期定常状態\(i_{ss}^H\)へ回帰する。
            \item[\(\hat{a}_1^H\)が負になることによるYS-\(\hat{p}\),CS-\(\hat{p}\)曲線の上移動]
            \label{itm:chapter5_a_shock_interpretations_rise_first_a_H_YS_CS}
                \(\hat{a}_1^H\)が負になることを受けて供給側では以下のような反応が起こる。
                \(\hat{p}_0^H=0\)であり、また\(\hat{a}_1^H\)が動き始めるこの時点においてまだ
                \(\operatorname{E}_1[\hat{\pi}_2^H]=0\)である。
                したがってYS-\(\hat{p}\),CS-\(\hat{p}\)曲線の切片の主決定要因
                \(\hat{p}_0^H+\beta_{ss}^H\operatorname{E}_1[\hat{\pi}_2^H]-2\kappa^H\hat{a}_1^H
                =-2\kappa^H\hat{a}_1^H\)
                が正となる。
                よってYS-\(\hat{p}\),CS-\(\hat{p}\)曲線は切片が上昇し、上へ移動する。
                これにより\(y_1^H,c_1^{H\to W}\)は下落し、\(p_1^H,p_1^{H\to W}\)は上昇する。
            \item[\(\hat{p}_1^H\)が正になることによるYS-\(\hat{p}\),CS-\(\hat{p}\)曲線のさらなる上移動]
            \label{itm:chapter5_a_shock_interpretations_rise_first_p_H_YS_CS}
                \(\hat{\pi}_1^H=\hat{p}_1^H-\hat{p}_0^H=\hat{p}_1^H\)が正となった。
                そこで\(\operatorname{E}_1[\hat{\pi}_2^H]\)も正であると仮定する。
                するとYS-\(\hat{p}\),CS-\(\hat{p}\)曲線の切片の主決定要因
                \(\hat{p}_0^H+\beta_{ss}^H\operatorname{E}_1[\hat{\pi}_2^H]-2\kappa^H\hat{a}_1^H
                =\beta_{ss}^H\operatorname{E}_1[\hat{\pi}_2^H]-2\kappa^H\hat{a}_1^H
                \approx\operatorname{E}_1[\hat{\pi}_2^H]-2\kappa^H\hat{a}_1^H\)
                がさらに上昇する。
                よってYS-\(\hat{p}\),CS-\(\hat{p}\)曲線は切片がさらに上昇し、さらに上へ移動する。
                これにより\(y_1^H,c_1^{H\to W}\)はさらに減少し、\(p_1^H,p_1^{H\to W}\)はさらに増加する。
            \item[\(i_1^H\)の下落したことによるYD,CD曲線の上移動圧力]
            \label{itm:chapter5_a_shock_interpretations_fall_first_i_H_YD_CD}
                このとき需要側においても変化が生じる。
                \(i_1^H\)が下落するため、YD,CD曲線の係数の主決定要因
                \(\frac{1}{(1+i_1^H)\operatorname{E}_1[\lambda_2^H]}\)
                において分母が減少し、係数が上昇する圧力が生じる。
                これはYD,CD曲線を上へ移動させる要因となる。
            \item[\(\operatorname{E}_1[\lambda_2^H]\)の大幅な上昇によるYD,CD曲線の強烈な下移動]
            \label{itm:chapter5_a_shock_interpretations_fall_first_E_lambda_YD_CD}
                しかし、これらの政策は将来の強力な金融引き締めと、
                それに伴う極端なデフレ（物価の大幅な下落）を予想させる。
                自国コンティンジェント債券に関するオイラー方程式
                \eqref{form:chapter5_beta_shock_policies_overview_euler_contingent_lim}
                において、将来の物価水準が極めて低くなるという予想から、
                期待限界効用 \(\operatorname{E}_1[\lambda_2^H]\) が大幅に上昇する。
                この効果が \(i_1^H\) の下落効果を完全に凌駕するため、
                YD,CD曲線の係数は結果として大幅に下落する。
                よってYD,CD曲線は極めて強く下へ移動する。
                これにより\(y_1^H,c_1^{H\to W}\)および\(p_1^H,p_1^{H\to W}\)に強烈な下落圧力が生じる。
            \item[1期のまとめ]
            \label{itm:chapter5_a_shock_interpretations_fall_first_summary}
                YS-\(\hat{p}\),CS-\(\hat{p}\)曲線の上移動（上昇圧力）と、
                YD,CD曲線の強烈な下移動（下落圧力）が同時に発生する。
                これらを合わせることにより、
                数量 \(y_1^H, c_1^{H\to W}\) については、供給と需要の双方からの下落圧力が重なり合い、
                極端なマイナス方向への大暴落を引き起こす。
                一方で価格 \(p_1^H, p_1^{H\to W}\) については、供給側の上昇圧力と需要側の下落圧力が拮抗するが、
                初期の供給制約の圧力が勝り、一時的にプラス方向へスパイクする。
        \end{description}

    \item
    \label{itm:chapter5_a_shock_interpretations_fall_second}
        \textbf{2期から谷の底まで：デフレの顕在化と実体経済の低迷}
        \begin{description}[style=nextline]
            \item[\(\hat{y}_t^H\)が極端な負になったことによるYS-\(\hat{p}\),CS-\(\hat{p}\)曲線の下移動]
            \label{itm:chapter5_a_shock_interpretations_fall_second_y_H_YS_CS}
                1期の大暴落により \(\hat{y}_t^H\) は極端な負の値となる。
                この期間においてYS-\(\pi\)曲線（差分形式）の右辺
                \(\hat{y}_t^H-\hat{a}_t^H\) は強い負の値となる。
                これにより \(\operatorname{E}_t[\hat{\pi}_{t+1}^H]\) が強烈な負（デフレ期待）となり、
                YS-\(\hat{p}\),CS-\(\hat{p}\)曲線の切片を押し下げる。
                よってYS-\(\hat{p}\),CS-\(\hat{p}\)曲線は下へ移動を開始する。
                これにより価格 \(p_t^H, p_t^{H\to W}\) は急反落し、デフレへと落ち込む。
            \item[\(\hat{a}_t^H\)の増加によるYS-\(\hat{p}\),CS-\(\hat{p}\)曲線のさらなる下移動]
            \label{itm:chapter5_a_shock_interpretations_fall_second_a_H_YS_CS}
                \(\hat{a}_t^H\)が負から\(0\)に向かって増加していくため、
                \(-2\kappa^H\hat{a}_t^H\) も減少していく。
                これによりYS-\(\hat{p}\),CS-\(\hat{p}\)曲線はさらに下へ移動する。
            \item[YD,CD曲線の下移動の継続]
            \label{itm:chapter5_a_shock_interpretations_fall_second_E_lambda_YD_CD}
                デフレへの突入と長期的な不況の予想により、
                期待限界効用 \(\operatorname{E}_t[\lambda_{t+1}^H]\) は依然として高い水準に留まる。
                そのため、YD,CD曲線は下移動した位置で圧力をかけ続ける。
            \item[2期から谷の底までのまとめ]
            \label{itm:chapter5_a_shock_interpretations_fall_second_summary}
                YS-\(\hat{p}\),CS-\(\hat{p}\)曲線が下移動に転じ、
                YD,CD曲線が低い位置に留まることにより、
                価格 \(p_t^H, p_t^{H\to W}\) は初期のスパイクから一転してデフレ領域へと急落する。
                数量 \(y_t^H, c_t^{H\to W}\) は、需要の低迷が続く中で深い谷の底に向かって
                さらに落ち込むか、低水準で推移し続ける。
        \end{description}

    \item
    \label{itm:chapter5_a_shock_interpretations_fall_third}
        \textbf{谷の底から収束まで：極めて緩やかな定常状態への回復}
        \begin{description}[style=nextline]
            \item[負の\(\hat{a}_t^H\)の\(0\)への接近と\(i_t^H\)の回復]
            \label{itm:chapter5_a_shock_interpretations_fall_third_a_H_i_H}
                負の\(\hat{a}_t^H\)は自己回帰過程にしたがい\(0\)へと接近していく。
                極端に低下していた \(i_t^H\) も、長い時間をかけて徐々に初期定常状態へ向かって回復していく。
            \item[\(\hat{y}_t^H-\hat{a}_t^H\)の縮小によるYS-\(\hat{p}\),CS-\(\hat{p}\)曲線の下移動の減衰]
            \label{itm:chapter5_a_shock_interpretations_fall_third_y_H_YS_CS}
                \(\hat{y}_t^H\) が谷の底を打ち、\(\hat{a}_t^H\) と共に徐々に回復していく過程で、
                負であった \(\hat{y}_t^H-\hat{a}_t^H\) の絶対値が縮小していく。
                これにより強烈であったデフレ期待 \(\operatorname{E}_t[\hat{\pi}_{t+1}^H]\) が減衰し、
                YS-\(\hat{p}\),CS-\(\hat{p}\)曲線の下移動のペースが落ちていく（停止へ向かう）。
            \item[\(\operatorname{E}_t[\lambda_{t+1}^H]\)の下落によるYD,CD曲線の上移動]
            \label{itm:chapter5_a_shock_interpretations_fall_third_E_lambda_YD_CD}
                将来のデフレ期待が徐々に剥落し、経済が定常状態へ向かうという予想が形成されるにつれ、
                異常に高止まりしていた期待限界効用 \(\operatorname{E}_t[\lambda_{t+1}^H]\) が下落していく。
                これによりYD,CD曲線の係数が上昇に転じ、YD,CD曲線は上へ移動を開始する。
            \item[谷の底から収束までのまとめ]
            \label{itm:chapter5_a_shock_interpretations_fall_third_summary}
                デフレ圧力の減衰に伴うYS-\(\hat{p}\),CS-\(\hat{p}\)曲線の下移動の停止と、
                期待の正常化に伴うYD,CD曲線の上移動が組み合わさる。
                これらにより、大暴落していた数量 \(y_t^H, c_t^{H\to W}\) は谷の底から反転し、
                極めて長い時間をかけて定常状態へと増加（回復）していく。
                価格 \(p_t^H, p_t^{H\to W}\) も深いマイナス圏から脱し、
                初期定常状態に向けて収束していく。
        \end{description}
\end{enumerate}